\section{Projektbeschreibung und Kontext}
\label{chap:projekt}

\subsection{Einleitung}

Im Rahmen unseres Projekts entwickeln wir eine Smartphone-Anwendung zur Objekterkennung. Als Anwendungsfeld haben wir die Hautkrebsvorsorge gewählt, konkret die Erkennung maligner Hautläsionen auf Fotos. SkinSentry wertet Aufnahmen von Hautveränderungen mit KI aus, schätzt das Hautkrebsrisiko in drei Stufen ein und empfiehlt bei Auffälligkeiten den Besuch bei einem Arzt.

Die Marktanalyse~\cite{marktanalyse2026} hat gezeigt, dass es bereits mehrere Apps zur Erkennung von Hautkrebs gibt (siehe \secref{abgrenzung-apps}). Aus dem Vergleich ihrer Funktionen haben wir eine grundlegende Entscheidung abgeleitet: Unsere App stellt keine Diagnose, sondern liefert eine Einschätzung und empfiehlt auf dieser Grundlage gegebenenfalls einen Arztbesuch. Für wen die App gedacht ist und welche Anforderungen sich daraus ergeben, beschreibt \chapref{zielgruppe}.

\subsection{Vision}

Langfristig soll die App ein digitaler Begleiter für die Hautkrebsvorsorge werden. Nutzer dokumentieren ihre Muttermale regelmäßig, Veränderungen werden automatisch erkannt, und die Ergebnisse lassen sich als Bericht an Hautärzte weitergeben. Über den Kurs hinaus wären eine Anbindung an Teledermatologie-Dienste sowie eine klinische Validierung und Zulassung als Medizinprodukt denkbar.

\subsection{Nutzungskontext}

Die App wird vor allem im Alltag zu Hause genutzt, etwa bei der regelmäßigen Selbstuntersuchung der Haut alle drei Monate. Zusätzlich kann sie in der Freizeit hilfreich sein, zum Beispiel im Sommer oder im Urlaub bei hoher UV-Belastung. Die App richtet sich an Erwachsene jeden Alters, ältere Menschen ausdrücklich eingeschlossen (siehe \chapref{zielgruppe}). Deshalb legen wir Wert auf eine einfache Bedienung: große Schaltflächen, eine klare Sprache und eine Schritt-für-Schritt-Anleitung bei der Fotoaufnahme.

\subsection{Ähnliche Apps und Abgrenzung}
\label{sec:abgrenzung-apps}

Zu den bekanntesten Anwendungen gehören SkinVision~\cite{app-skinvision} und SkinScreener~\cite{app-skinscreener}, die eine KI-basierte Risikoeinschätzung anbieten, sowie Miiskin~\cite{app-miiskin}, das sich auf die Dokumentation von Muttermalen konzentriert. Teledermatologie-Dienste wie Dermanostic~\cite{app-dermanostic} und OnlineDoctor~\cite{app-onlinedoctor} verbinden Nutzer dagegen direkt mit Hautärzten. Unsere App unterscheidet sich vor allem durch die Verbindung von Risikoeinschätzung und Verlaufsbeobachtung in einer einfachen Oberfläche, die auch ältere Nutzer ohne Vorkenntnisse bedienen können. Außerdem bleiben alle Daten auf dem Smartphone: Die KI-Analyse läuft direkt auf dem Gerät, Fotos und Verlauf werden nur lokal gespeichert, und ein Nutzerkonto ist nicht nötig (siehe \chapref{architektur}).
